\documentclass[11pt]{article}
\usepackage[margin=25mm]{geometry}
\usepackage{fontspec}
\setmainfont{Latin Modern Roman}
\usepackage{amsmath,amssymb,mathtools}
\usepackage{unicode-math}
\setmathfont{Latin Modern Math}
\usepackage{xurl,hyperref}
\hypersetup{colorlinks=true,urlcolor=blue}
\setcounter{secnumdepth}{0}
\setlength{\parindent}{0pt}
\setlength{\parskip}{5pt}
\setlength{\emergencystretch}{3em}
\newcommand{\sref}[2]{\hyperlink{source:#1:#2}{[S#2]}}
\newcommand{\cataloguescope}[1]{\paragraph{Recorded target.}#1}
\newcommand{\R}{\mathbb R}
\newcommand{\C}{\mathbb C}
\newcommand{\N}{\mathbb N}
\newcommand{\Z}{\mathbb Z}
\newcommand{\Q}{\mathbb Q}
\newcommand{\F}{\mathbb F}
\newcommand{\eps}{\varepsilon}
\begin{document}
% UnsolvedMath ID 30004443; source catalogue code OWR-17476-003
\setcounter{section}{30004442}
\section{Numerical Schemes for Coupled Fracture and Dissipation}\label{30004443}
{\small UnsolvedMath ID 30004443 · Mathematical Physics}
\subsection{Definitions and mathematical statement}

\paragraph{Shared notation.}$\mathbb N=\{1,2,\ldots\}$; the standard number fields have their usual meanings. All additional conventions are stated locally.


\paragraph{Definitions and conventions.}
A rate-independent model has dissipation positively homogeneous of degree one in the evolution rate, so its response is invariant under suitable time reparametrization. Rate-dependent processes have a distinct time scale, often through viscous dissipation. Fracture and damage variables obey irreversibility constraints. Energetic, local and vanishing-viscosity solution concepts may differ at jumps when the energy is nonconvex. A reliable numerical scheme must identify and converge to the intended concept while resolving the coupling.

\paragraph{Target retained from the canonical source catalogue.}
\paragraph{Catalogue wording retained for traceability.}
For rate-independent fracture and damage models coupled with rate-dependent processes, are there efficient feasible numerical schemes that reliably approximate the intended solution class, and what mathematical and mechanical framework should describe the coupling?

\paragraph{Source scope and qualification.}
The desired solution class must be selected explicitly: the report emphasizes that energetic and vanishing-viscosity evolutions can have different jumps. Efficiency alone is not convergence to the intended physical and mathematical evolution.
\subsection{Short English statement}
Develop a sound coupled framework and practical convergent schemes for fracture or damage interacting with time-dependent dissipation.
\subsection{Research attempt}
\paragraph{Model, date and investigated scope.}
OpenAI GPT-6 Astra (configured model \texttt{gpt-6-astra}), ultra reasoning; generated 11 October 2026. Derived an exact incremental energy/work estimate for a simple irreversible-plus-viscous scheme; no convergence to the intended fracture jump concept.

\paragraph{Sources and literature check.}
Read the supplied definition, catalogue wording, source qualifications and reference entries for 30004443. The historical source/context and inherited status are retained; no fresh literature or novelty certification is claimed. No external theorem is used without proof in this attempt; explicit model assumptions are distinguished from facts about the source model.

\paragraph{A coupled incremental energy inequality.} For scalar displacement \(u\), irreversible damage variable \(z\), step \(h>0\), viscosity \(\nu>0\), toughness \(r\ge0\), and load \(\ell_n\), take
\[
 E_n(u,z)=\tfrac12(u-z)^2+W(z)-\ell_nu.
\]
Assume a global minimizer \((u_n,z_n)\) exists for
\[
 E_n(u,z)+r(z-z_{n-1})+\frac{\nu}{2h}(u-u_{n-1})^2,
 \qquad z\ge z_{n-1}.
\]
The previous pair is an admissible competitor. Therefore, with \(\Delta u_n=u_n-u_{n-1}\) and \(\Delta z_n=z_n-z_{n-1}\),
\[
 E_n(u_n,z_n)+r\Delta z_n+\frac\nu{2h}|\Delta u_n|^2
 \le E_{n-1}(u_{n-1},z_{n-1})-(\ell_n-\ell_{n-1})u_{n-1}.
\]
Summation telescopes the energies and displays the discrete work contribution explicitly. No differentiability of \(W\) is needed for this comparison.
\paragraph{Why the limit is not determined.} This estimate is compatible with global-minimum jumps in \(z\). If the intended evolution selects local metastable branches or a vanishing-viscosity jump cost, the inequality does not show that this globally minimizing update converges to it. Different ratios of numerical step to regularizing viscosity can retain different transition behavior.
\paragraph{Scope audit.} The model is a scalar coupled test, not a fracture-space discretization or an efficient algorithm. Minimizer existence, mesh compactness and limiting energy/work identities require additional coercivity and convergence proofs. The first gap is choosing the source solution concept and proving the time/space/viscosity limits select precisely it. The derived inequality alone cannot certify reliability or mechanical validity.

\paragraph{Verification and adversarial audit.}
Checked admissibility of the old state, the sign of load work, telescoping and the conditional nature of minimizer existence. This is a GPT-6 Astra Ultra self-audit, not independent human verification. Prompt A drafting followed by prompt B scope/dependency/boundary self-audit. Checked admissibility of the old state, the sign of load work, telescoping and the conditional nature of minimizer existence. The remaining gap is: Prove the coupled discretization and viscosity limits select the specified energetic/local/vanishing-viscosity solution class. Unresolved status and zero original-target completion are retained.

Coordinator checked the competitor inequality and sign of discrete work, including its dependence on the assumed minimizer. This additional check was model-assisted, not independent human verification.

\paragraph{Tools, process and contributions.}
Python 3 and local file-reading/writing tools for manuscript serialization and structural checks; mathematical calculations were symbolic reasoning, with no numerical experiment or theorem prover. Contribution prompt A was followed by prompt B. GPT-6 Astra with ultra reasoning drafted the research and performed its stated self-audit; any additional coordinator contribution is identified in the verification and audit. A GPT-6-based Codex coordinator assembled the manuscript and checked repository integration; its exact variant and reasoning configuration were not exposed. No human mathematical contribution, independent human proof check or formal proof certification is supplied. Complete derivations and the consulted source references are included here; private reasoning transcripts are not published.

\paragraph{Final status and first remaining gap.}
\textbf{UNRESOLVED}. Prove the coupled discretization and viscosity limits select the specified energetic/local/vanishing-viscosity solution class.

\section{Completion Estimate}
\noindent\textbf{COMPLETION: 0\%}
The proved calculation is an elementary auxiliary result or a test of a restricted strategy; no new portion of the full source target is established.

\subsection{Sources}
\begin{itemize}
\item[\hypertarget{source:30004443:1}{S1}] ULAM.AI and the UnsolvedMath Contributors, UnsolvedMath, record 30004443 (OWR-17476-003). Catalogue curation under CC BY 4.0; definitions and English exposition added locally.
\url{https://huggingface.co/datasets/ulamai/UnsolvedMath}.
\item[\hypertarget{source:30004443:2}{S2}] Mechanics of Materials: Towards Predictive Methods for Kinetics in Plasticity, Fracture, and Damage (2020), Oberwolfach Reports, DOI 10.4171/owr/2020/13, pp. 713--783. Primary report accessed 2026-10-11; relevant PDF pages 35.
\url{https://doi.org/10.4171/owr/2020/13}.
\item[\hypertarget{source:30004443:3}{S3}] Rossi--Thomas (2017) (catalogue-listed context; not independently consulted in this preparation).
\url{https://doi.org/10.1137/15M1051567}.
\end{itemize}
\label{end:30004443}
\end{document}
